\appendix
\section*{Phụ lục}

\section{Mã nguồn trọng tâm (rút gọn)}

\subsection{\texttt{rgb2gray.v}}
\begin{lstlisting}[style=vncode]
// Y ≈ (77*R + 150*G + 29*B) >> 8
module rgb2gray(
  input  [7:0] r0,g0,b0, r1,g1,b1,
  output [7:0] R0,G0,B0, R1,G1,B1
);
  integer y0,y1;
  function [7:0] clamp8(input integer v);
    begin if(v<0) clamp8=0; else if(v>255) clamp8=255; else clamp8=v[7:0]; end
  endfunction
  always @* begin
    y0 = (77*r0 + 150*g0 + 29*b0) >> 8;
    y1 = (77*r1 + 150*g1 + 29*b1) >> 8;
  end
  assign R0=clamp8(y0); assign G0=clamp8(y0); assign B0=clamp8(y0);
  assign R1=clamp8(y1); assign G1=clamp8(y1); assign B1=clamp8(y1);
endmodule
\end{lstlisting}

\subsection{\texttt{rgb2ycbcr.v}}
\begin{lstlisting}[style=vncode]
// BT.601 limited-range (fixed-point 1/256)
module rgb2ycbcr(
  input  [7:0] r0,g0,b0, r1,g1,b1,
  output [7:0] Y0,Cb0,Cr0, Y1,Cb1,Cr1
);
  integer y0,cb0,cr0,y1,cb1,cr1;
  function [7:0] clamp8(input integer v);
    begin if(v<0) clamp8=0; else if(v>255) clamp8=255; else clamp8=v[7:0]; end
  endfunction
  always @* begin
    y0  = ( 77*r0 + 150*g0 + 29*b0) >>> 8;
    y1  = ( 77*r1 + 150*g1 + 29*b1) >>> 8;
    cb0 = ((-43*r0 - 85*g0 + 128*b0) >>> 8) + 128;
    cb1 = ((-43*r1 - 85*g1 + 128*b1) >>> 8) + 128;
    cr0 = ((128*r0 - 107*g0 -  21*b0) >>> 8) + 128;
    cr1 = ((128*r1 - 107*g1 -  21*b1) >>> 8) + 128;
  end
  assign Y0=clamp8(y0);  assign Cb0=clamp8(cb0); assign Cr0=clamp8(cr0);
  assign Y1=clamp8(y1);  assign Cb1=clamp8(cb1); assign Cr1=clamp8(cr1);
endmodule
\end{lstlisting}

\section{Header Format BMP 24-bit}
\begin{table}[H]
\centering
\renewcommand{\arraystretch}{1.3}
\begin{tabular}{@{}c c c p{0.45\linewidth}@{}}
\toprule
\textbf{Offset} & \textbf{Kích thước (byte)} & \textbf{Trường} & \textbf{Mô tả} \\
\midrule
0  & 2 & \texttt{Signature} & Chữ ký định dạng, luôn là \texttt{4D42} (``BM'') \\
2  & 4 & \texttt{FileSize} & Kích thước toàn bộ tệp BMP (có thể không chính xác) \\
6  & 2 & \texttt{Reserved1} & Dự phòng, phải bằng 0 \\
8  & 2 & \texttt{Reserved2} & Dự phòng, phải bằng 0 \\
10 & 4 & \texttt{DataOffset} & Vị trí bắt đầu vùng dữ liệu ảnh (byte) \\
14 & 4 & \texttt{DIBSize} & Kích thước cấu trúc \texttt{BITMAPINFOHEADER}, luôn là 40 \\
18 & 4 & \texttt{Width} & Chiều rộng ảnh (pixel) \\
22 & 4 & \texttt{Height} & Chiều cao ảnh (pixel) \\
26 & 2 & \texttt{Planes} & Số mặt phẳng ảnh, luôn là 1 \\
28 & 2 & \texttt{BitsPerPixel} & Số bit trên mỗi pixel (1, 4, 8 hoặc 24) \\
30 & 4 & \texttt{Compression} & Loại nén (0=không nén, 1=RLE-8, 2=RLE-4) \\
34 & 4 & \texttt{ImageSize} & Kích thước vùng dữ liệu ảnh (gồm padding) \\
38 & 4 & \texttt{XResolution} & Độ phân giải ngang (pixel/mét), có thể không chính xác \\
42 & 4 & \texttt{YResolution} & Độ phân giải dọc (pixel/mét), có thể không chính xác \\
46 & 4 & \texttt{ColorsUsed} & Số màu sử dụng, hoặc 0 nếu dùng đầy đủ \\
50 & 4 & \texttt{ImportantColors} & Số màu quan trọng, hoặc 0 \\
\bottomrule
\end{tabular}
\caption{Cấu trúc Header BMP 24-bit (File Header + DIB Header)}
\label{tab:bmp_header}
\end{table}

\section{Quy trình tái lập}
\begin{enumerate}
  \item Cài Python \& Pillow: \texttt{pip install pillow}.
  \item Chạy \texttt{image\_to\_hex.py}; chọn ảnh cần chuyển đổi, chọn kích thước ảnh muốn chuyển, bật RESIZE nếu cần; lưu file với tên \texttt{input.hex} tại thư mục \texttt{img/}.
  \item Mở Vivado, import các file .v trong \texttt{rtl/}, \texttt{rgb_convert/}, \texttt{sim/}; đặt \texttt{tb\_simulation} là top của \texttt{sim\_1}.
  \item Chỉnh \texttt{parameter.v} (đường dẫn \& kích thước), bật macro xử lý mong muốn.
  \item Run simulation $\approx$ 6\,ms hoặc -all; đóng simulator sau khi thấy \texttt{write\_done\_*} lên 1.
  \item Kiểm tra các file \texttt{output\_*.bmp}.
\end{enumerate}

\section{Ghi chú/Checklist}
\begin{itemize}
  \item \textbf{Clamp:} luôn giới hạn 0--255 để tránh tràn khi quy đổi.
  \item \textbf{Padding:} mỗi dòng BMP phải là bội số của 4~byte.
  \item \textbf{Kích thước:} giữ đồng bộ giữa script, \texttt{parameter.v} và \texttt{image\_write.v} (Đã tham số hóa).
  \item \textbf{Đường dẫn:} ưu tiên tương đối; tránh hard-code tuyệt đối trong testbench.
\end{itemize}

\addcontentsline{toc}{section}{Nguồn tham khảo}

\nocite{*}  % <-- dòng này giúp in toàn bộ tài liệu trong file .bib

\bibliographystyle{ieeetr}
\bibliography{ref.bib}